\scanpage{13}
$\mathrm{NS}(\widetilde X)'\longrightarrow J'(k(\eta))$ est égal au groupe engendré
par l’image de $\tilde\xi_0$, car si $\alpha$ est
dans le noyau, alors $\alpha$ provient de
$\alpha_0\in\operatorname{Pic}(\widetilde X)'$, et comme p.h. l’image de
$\alpha_0$ dans $\deleted{\unsure{\operatorname{Pic}}}J(Y_\eta)$ est dans $J(Y_\eta)^{\rm fixe}\simeq\operatorname{Pic}^0(\widetilde X)$,
il s’ensuit que quitte à modifier $\alpha_0$ par
un \unsure{élmt|élément} de $\operatorname{Pic}^0(\widetilde X)$, on aura une
image nulle \ldots{} Donc sous réserve
que tous les $Y_y$ soient intègres, le
noyau de $\mathrm{NS}(\widetilde X)'\unsure{\text{---}|\longrightarrow}J'(k(\eta))$ est
de rang $1$, et engendré par l’image de $\bar\xi_0$.

Enfin, je dis que \deleted{l’image de} $\mathrm{NS}(\widetilde X)'\longrightarrow J'(k(\eta))$
est \uline{surjectif} \inserted{mod groupe fini}. En effet, \deleted{un élément de
$J'(k(\eta))$} comme $H^1(K,J'')$ est un
groupe de torsion \unsure{[|$T$} où $K=k(\eta)=k(t)$,
$J''=J^{\rm fixe}$) on voit que mod groupes finis,
\deleted{$\mathrm{NS}(\widetilde X)'\text{---}J'(k(\eta))$ \unsure{v|est}} $J(k(\eta))\text{---}J'(k(\eta))$ est
surjectif, donc il suffit de prouver
que $\operatorname{Pic}(\widetilde X)'\longrightarrow J(k(\eta))$ est surjectif \deleted{\unsure{mod}}
% archive sheet number: 12
